\documentclass[a4paper, titlepage]{article}

\usepackage[utf8x]{inputenc} % accents
\usepackage[T1]{fontenc}      % caractères français
\usepackage[margin=1in]{geometry}         % marges
\usepackage[francais]{babel}  % langue
\usepackage{graphicx,subfigure}         % images
\usepackage{verbatim}         % texte préformaté
\usepackage{float}
\usepackage{framed}	% texte encadré

\usepackage{fancyhdr}
\pagestyle{fancy}
\usepackage{lastpage}

\renewcommand\headrulewidth{1pt}
\fancyhead[L]{Algorithmique avancée -- Ordonnancement simple}
\fancyhead[R]{\bsc{Enssat}}

\renewcommand\footrulewidth{1pt}
\fancyfoot[C]{\today}
\fancyfoot[L]{\bsc{Nokaya} -- \bsc{Vythelingum}}
\fancyfoot[R]{\textbf{Page \thepage/\pageref{LastPage}}}

\title{Ordonnancement simple}
\author{Kévin \bsc{Vythelingum}, Jean-Michel \bsc{Nokaya}}
\date{\today}           % la future date de parution



\begin{document}
\makeatletter
  \begin{titlepage}
  \centering
  	{\large \textsc{École Nationale Supérieure des Sciences Appliquées et de Technologie}}\\
	\textsc{Logiciel et Systèmes Informatiques - Deuxième année}\\
  \vfill
	\textbf{Algorithmique avancée}\\
  \vspace{1cm}
        {\LARGE \textbf{\@title}} \\
  \vspace{2em}
  {\large Kévin \bsc{Vythelingum}}\\
  \vspace{0.5cm}
  {\large Jean-Michel \bsc{Nokaya}} \\
  \vspace{1cm}

	\@date \\
  \vfill
	\includegraphics[height=0.07\textheight]{enssat.png}
	\hfill
	\begin{tabular}{l}
		\large Chargé de cours :\\[0.2cm]
		\large Olivier \bsc{Pivert} \\
		\vspace{1cm}
	\end{tabular}

\end{titlepage}\makeatother
\tableofcontents
\newpage
\large

\section{Introduction}

Ce projet consiste à examiner deux manières différentes pour résoudre le même problème afin de comparer leurs complexités.
Ainsi, nous nous intéressons à un problème d'ordonnancement simple, où différentes tâches peuvent être exécutées avec des durées différentes, moyennant un coût plus important pour une durée plus courte.
Il s'agit de déterminer le meilleur ordonnancement pour avoir un coût minimal en s'étant fixé une durée maximale d'exécution.
Pour cela, nous réaliserons un premier algorithme suivant la méthode des essais successifs, puis nous élaborerons un algorithme de programmation dynamique.

\section{Préliminaires}

\begin{enumerate}
\item
	On cherche tout d'abord à montrer que la solution optimale recherchée est trouvée pour $\sum_{i=1}^{n} d_i = D$.

	Premièrement, il est important de remarquer qu'on a nécessairement $\sum_{i=1}^{n} d_i \le D$ car D est la durée globale maximale.
	Or, si on considère une solution S telle que $\sum_{i=1}^{n} d_i < D$, alors il existe au moins un $d_i$ qui ne soit pas maximal par rapport aux durées satisfaisantes pour la tâche $T_i$.
	Il existe donc une solution S' telle que $\sum_{i=1}^{n} d'_i \le D$ et $\sum_{i=1}^{n} d'_i > \sum_{i=1}^{n} d_i$.
	Comme le coût diminue lorsque la durée augmente, S' est donc une meilleure solution que S.
	S n'est donc pas optimale.

	On en déduit que la solution optimale recherchée est trouvée pour $\sum_{i=1}^{n} d_i = D$.

\item
	Ensuite, on s'intéresse à la complexité au pire d'une solution de type essais successifs en termes de nombre de candidats à examiner.

	La complexité au pire correspond à considérer toutes les combinaisons différentes possibles.
	Il s'agit donc de dénombrer le nombre N de candidats.

	Construire une solution consiste à choisir une durée pour $T_1$ parmi les k durées possibles, puis une durée pour $T_2$, et ainsi de suite jusq'à $T_n$.
	On a donc :

	\begin{eqnarray*}
		N & = & \underbrace{k \times k \times ... \times k}_{n fois} \\
		N & = & k^n
	\end{eqnarray*}

	La complexité recherchée vaut donc $k^{n}$.

\end{enumerate}

\section{Méthodes des essais successifs}

	\subsection{Analyse}

		On commence par définir les différents éléments qui interviennent dans le cadre de la méthode des essais successifs.
		En particulier, il s'agit de définir ce qu'est un vecteur représentant un candidat (solution) et quels sont les constituant de l'algorithme générique ($S_{i}, satisfaisant, enregistrer, soltrouvee, defaire$).

		\paragraph{}\noindent
		Solution : un candidat est un vecteur de taille n où chaque coefficient est une durée choisie parmi l'ensemble \{1,...,k\} (à chaque tâche on associe une durée).
		On choisit d'enregistrer les choix réalisés dans un tableau T de taille n.

		\paragraph{}\noindent
		$S_{i}$ : l'ensemble des durées possibles de 1 à k

		\paragraph{}\noindent
		$satisfaisant(x_{i}) = \sum_{j=1}^{i} x_{j} \le D$

		(la somme partielle des durées choisies est inférieure à la durée maximale autorisée)

		\paragraph{}\noindent
		$enregistrer(x_{i}) = T[i] \leftarrow x_{i}$

		\paragraph{}\noindent
		\emph{soltrouvée} : $i = n$

		\paragraph{}\noindent
		$defaire(x_{i}) = T[i] \leftarrow 0$

		\paragraph{}
		Pour simplifier les vérifications au niveau de satisfaisant et des conditions d'élagage, on utilisera les variables entières \emph{cout} et \emph{duree} initialisée à 0,
		qui représenteront le coût courant et la durée courante duent à nos choix de durées.
		Elles seront mises à jour dans \emph{enregistrer} et dans \emph{défaire}.
		En effet, en notant \emph{CD} le tableau à deux dimensions ayant les coûts en ligne et les durées en colonne, on effectura dans \emph{enregistrer} :

			\begin{eqnarray*}
				\mbox{coût}  & \leftarrow & \mbox{coût} + CD[i][x_{i}] \\
				\mbox{durée} & \leftarrow & \mbox{durée} + x_{i}
			\end{eqnarray*}

		Ensuite on effectura dans \emph{défaire} :

			\begin{eqnarray*}
				\mbox{coût}  & \leftarrow & \mbox{coût} - CD[i][x_{i}] \\
				\mbox{durée} & \leftarrow & \mbox{durée} - x_{i}
			\end{eqnarray*}


	\subsection{Condition d'élagage}
	On a montré dans les préliminaires que la solution optimale recherchée est trouvée pour $\sum_{i=1}^{n} d_i = D$.
	Cela va nous servir de base pour notre condition d'élagage.
	En effet, on va élaguer si on ne peut pas atteindre D avec les durées restantes à choisir.
	Autrement dit, en notant $S_1$ la somme des durées choisies et $S_2$ la somme des durées maximales restantes, on élaguera si $S_1 + S_2 < D$.
	Cependant, lorsque D est grand (i.e. qu'il est impossible de l'atteindre, même en choisissant toutes les durées maximales), il ne faudra pas élaguer.

	L'intérêt de cette condition d'élagage est d'arrêter plus tôt la recherche d'une solution en prévoyant qu'elle ne pourra pas être optimale.
	Nous vérifierons par des expérimentations que cette condition améliore effectivement la complexité temporelle en terme de nombre d'appels récursifs.

	\subsection{Algorithmes}

		\subsubsection{Sans élagage}

			Cet algorithme repose sur le modèle de programmation de recherche d'une solution optimale par essais successifs.
			Le principe est d'essayer toutes les combinaisons possibles tant qu'elles sont satisfaisantes, c'est-à-dire dans notre cas que la somme partielle des durées choisies est inférieure à la durée maximale autorisée D.
			Pour obtenir la solution optimale, on vérifie que le coût d'une solution obtenue est strictement inférieur au coût de la meilleure solution trouvée, auquel cas on l'affiche.
			La dernière solution affichée sera donc la solution optimale.

			L'appel de l'algorithme consiste à initialiser le coût optimal à une très grande valeur, puis à appeler la procédure au rang 1, c'est-à-dire commencer par choisir la durée de la première tâche $T_{1}$.
			On initialisera également les variables \emph{cout} et \emph{duree} à 0.

			\begin{tabbing}

			\hspace{1cm}\=\hspace{1cm}\=\hspace{1cm}\=\hspace{1cm}\=\hspace{1cm}\=\kill % défini la longueur d'une tabulation

			\textbf{procédure} $ordonnancement\_simple$ (ent i); \\
			\textbf{var} ent $x_{i}$; \\
			\textbf{début} \\
			\> \textbf{pour} $x_{i}$ \textbf{de} 1 \textbf{à} k \textbf{faire} \\
			\>\> \textbf{si} $duree+x_{i}\le D$ \textbf{alors} \\
			\>\>\> $duree \leftarrow duree + x_{i};$ \\
			\>\>\> $cout \leftarrow cout + cd[i][x_{i}];$ \\
			\>\>\> $T[i] \leftarrow x_{i};$ \\

			\>\>\> \textbf{si} i = n \textbf{alors} \\
			\>\>\>\> \textbf{si} $cout < cout\_opt$ \textbf{alors} \\
			\>\>\>\>\> $cout\_opt \leftarrow cout$; \\
			\>\>\>\>\> $affiche(T)$; \\
			\>\>\>\> \textbf{fsi}; \\

			\>\>\> \textbf{sinon} \\
			\>\>\>\> $ordonnancement\_simple(i+1)$; \\
			\>\>\> \textbf{fsi}; \\

			\>\>\> $T[i] \leftarrow 0;$ \\
			\>\>\> $cout \leftarrow cout - cd[i][x_{i}];$ \\
			\>\>\> $duree \leftarrow duree - x_{i};	$ \\
			\>\> \textbf{fsi}; \\
			\> \textbf{fait}; \\
			\textbf{fin};

			\end{tabbing}

			\noindent
			Appel :\\
			$coutOpt \leftarrow \infty;$ \\
			$cout \leftarrow 0;$ \\
			$duree \leftarrow 0;$ \\
			$ordonnancement\_simple(1)$;

		\subsubsection{Avec élagage}

			Pour réaliser l'algorithme avec élagage, on dispose d'une fonction \emph{encorepossible(di,ci)} qui prend en argument la durée et le coût sélectionnés et retourne un booléen :
			\texttt{vrai} si on peut encore espérer trouver une solution optimale, \texttt{faux} s'il est nécessaire d'élaguer.
			En partant de l'algorithme précédent, effectue l'appel à \emph{ordonnancement\_simple} pour le rang suivant à condition que \emph{encorepossible} soit \texttt{vrai}.
			On obtient alors l'algorithme :

			\begin{tabbing}

			\hspace{1cm}\=\hspace{1cm}\=\hspace{1cm}\=\hspace{1cm}\=\hspace{1cm}\=\kill % défini la longueur d'une tabulation

			\textbf{procédure} $ordonnancement\_simple$ (ent i); \\
			\textbf{var} ent $x_{i}$; \\
			\textbf{début} \\
			\> \textbf{pour} $x_{i}$ \textbf{de} 1 \textbf{à} k \textbf{faire} \\
			\>\> \textbf{si} $duree+x_{i}\le D$ \textbf{alors} \\
			\>\>\> $duree \leftarrow duree + x_{i};$ \\
			\>\>\> $cout \leftarrow cout + cd[i][x_{i}];$ \\
			\>\>\> $T[i] \leftarrow x_{i};$ \\

			\>\>\> \textbf{si} i = n \textbf{alors} \\
			\>\>\>\> \textbf{si} $cout < cout\_opt$ \textbf{alors} \\
			\>\>\>\>\> $cout\_opt \leftarrow cout$; \\
			\>\>\>\>\> $affiche(T)$; \\
			\>\>\>\> \textbf{fsi}; \\

			\>\>\> \textbf{sinon} \\
			\>\>\>\> \textbf{si} $encorepossible (duree)$ \textbf{alors} \\
			\>\>\>\>\> $ordonnancement\_simple(i+1)$; \\
			\>\>\>\> \textbf{fsi}; \\
			\>\>\> \textbf{fsi}; \\

			\>\>\> $T[i] \leftarrow 0;$ \\
			\>\>\> $cout \leftarrow cout - cd[i][x_{i}];$ \\
			\>\>\> $duree \leftarrow duree - x_{i};	$ \\
			\>\> \textbf{fsi}; \\
			\> \textbf{fait}; \\
			\textbf{fin};

			\end{tabbing}

			\noindent
			Appel :\\
			$coutOpt \leftarrow \infty;$ \\
			$cout \leftarrow 0;$ \\
			$duree \leftarrow 0;$ \\
			$ordonnancement\_simple(1)$;

			\paragraph{}
			Le comportement de \emph{encorepossible} reprend le principe d'élagage énoncé dans la section précédente.
			On suppose que l'on dispose d'un tableau \emph{dmax} de taille n qui est initialisé à la durée maximale qu'il est possible d'allouer à chaque tâche,
			et d'une fonction \emph{sommeTableau(tableau,taille)} qui permet de calculer la somme des éléments d'un tableau d'entiers.

			\begin{tabbing}

			\hspace{1cm}\=\hspace{1cm}\=\hspace{1cm}\=\hspace{1cm}\=\hspace{1cm}\=\kill % défini la longueur d'une tabulation

			\textbf{fonction} $encorepossible$ (ent di, ent ci) : booléen; \\
			\textbf{var} ent S; \\
			\textbf{début} \\
			\> \textbf{si} (S < D) \textbf{et} (ci = 1) \textbf{alors} \\
			\>\> \textbf{retourner} \texttt{vrai}; \\
			\> \textbf{sinon} \\
			\>\> \textbf{si} $S-dmax[ci]+di < D$ \textbf{alors} \\
			\>\>\> \textbf{retourner} \texttt{faux}; \\
			\>\> \textbf{sinon} \\
			\>\>\> $dmax[ci] \leftarrow di;$ \\
			\>\>\> \textbf{retourner} \texttt{vrai}; \\
			\textbf{fin}

			\end{tabbing}

	\subsection{Expérimentations}

	Pour tester la validité de nos algorithmes et mesurer leurs performances, nous affichons la solution optimale proposée ainsi que le nombre d'appels effectués à la procédure \emph{ordonnancement\_simple}.
	Cela nous permettra de comparer les algorithmes avec et sans élagage.

	Tout d'abord, nous avons testé l'exemple proposé par le sujet, c'est-à-dire avec $n=4$, $k=5$, $D=10$, et le tableau suivant :

	\begin{center}
		\begin{tabular}{|c|c|c|c|c|c|}
			\hline
			cd	& 1	& 2	& 3	& 4	& 5 \\ \hline
			$c_1$	& 110	& 90	& 65	& 55	& -- \\ \hline
			$c_2$	& 120	& 90	& 70	& 50	& 40 \\ \hline
			$c_3$	& 90	& 70	& 65	& 60	& -- \\ \hline
			$c_4$	& 65	& 60	& 55	& --	& -- \\
			\hline
		\end{tabular}
	\end{center}

	Nous obtenons les résultats suivant pour les algorithmes sans et avec élagage :

\begin{framed}
\paragraph{Sans élagage\\}
\begin{verbatim}
3 4 2 1
duree = 10 cout = 250

Nombre d'appels : 121
\end{verbatim}
\end{framed}

\begin{framed}
\paragraph{Avec élagage\\}
\begin{verbatim}
3 4 2 1
duree = 10 cout = 250

Nombre d'appels : 99
\end{verbatim}
\end{framed}

	\paragraph{}
	Nous obtenons heureusement le même résultat avec les deux algorithmes, c'est-à-dire affecter respectivement les durées 3, 4, 2, 1 aux tâches $c_1$, $c_2$, $c_3$, $c_4$.
	De plus, nous arrivons à faire 1.2 fois moins d'appels avec la condition d'élagage, ce qui montre son intérêt.

	Nous effectuons ensuite des tests avec des jeux de données différents en vérifiant à chaque fois la cohérence des résultats obtenus.
	Pour cela, nous avons conçu un programme \emph{es-experimentations} qui génère des tableaux \emph{cd} avec des valeurs aléatoires comprises entre 10 et 1000 ou $+\infty$ et les applique sur les deux algorithmes.
	L'objectif est de comparer les performances des deux algorithmes.
	Ainsi, nous avons recensé les nombres d'appels obtenus pour des valeurs de k et n différentes dans le tableau suivant.

	\begin{center}
		\begin{tabular}{|c|c|c|c|r|r|}
			\hline
			Test	& n	& k	& D	& Nombre d'appels sans élagage & Nombre d'appels avec élagage \\
			\hline
			0 	& 4	& 5	& 10	& 121 	& 99 \\
			\hline
			1 	& 4	& 5	& 10 	& 121	& 113 \\
			2 	& 4	& 5	& 10 	& 121	& 51 \\
			3 	& 4	& 5	& 10 	& 121	& 101 \\
			4 	& 4	& 5	& 10 	& 121	& 113 \\
			5 	& 4	& 5	& 10 	& 121	& 117 \\
			\hline
			6 	& 10	& 5	& 30 	& 2 015 376 	& 274 900 \\
			7 	& 10	& 5	& 30 	& 2 015 376 	& 809 160 \\
			8 	& 10	& 5	& 30 	& 2 015 376 	& 275 006 \\
			\hline
			9 	& 10	& 10	& 30 	& 19 580 242 	& 8 151 656 \\
			10 	& 10	& 10	& 30 	& 19 580 242 	& 10 869 695 \\
			11 	& 10	& 10	& 30 	& 19 580 242 	& 372 187 \\
			\hline
			12 	& 10	& 20	& 30 	& 22 958 977 	& 394 759 \\
			13 	& 10	& 30	& 30 	& 22 964 087 	& 369 466 \\
			14 	& 10	& 50	& 30 	& 22 964 087 	& 369 246 \\
			\hline
			15 	& 10	& 100	& 30 	& 22 964 087 	& 931 513 \\
			16 	& 10	& 5000 	& 30 	& 22 964 087 	& 368 336 \\

			\hline
		\end{tabular}
	\end{center}
	
	On remarque deux choses.
	Tout d'abord, le nombre d'appels sans élagage ne dépend pas du jeu de données \emph{cd} étudié.
	Ensuite, un seuil est atteint pour l'algorithme sans élagage avec un grand nombre de durées possibles.
	Ceci peut être explqué par le fait que fixer D impose une certaine plage de valeurs possibles pour les durées : cela ne sert à rien de choisir comme première valeur un di supérieur à D, la solution ne pourrait alors pas être satisfaisante.
	On obtient donc le même nombre d'appels pour les cas où $k \ge D$.
	
	Il s'agit maintenant de fixer k à 10 et de faire varier n.

	\begin{center}
		\begin{tabular}{|c|c|c|c|r|r|}
			\hline
			Test	& n	& k	& D	& Nombre d'appels sans élagage & Nombre d'appels avec élagage \\
			\hline
			0 	& 5	& 10	& 30	& 10 396 	& 11 \\ \hline
			1 	& 10	& 10	& 30	& 19 580 242 	& 2 721 266 \\ \hline
			2 	& 15	& 10	& 30	& 449 171 192 	& 182 779 801 \\ \hline
			3 	& 20	& 10	& 30	& 1 010 258 502 	& 581 312 322 \\ \hline
			4 	& 50	& 10	& 30	& 1 063 267 584 	& 1 063 267 584 \\
			\hline
		\end{tabular}
	\end{center}
	
	On remarque qu'il n'est possible de monter que jusqu'à $n = 50$ dans un temps raisonnable.
	En effet, la complexité étant en $k^n$ d'après les questions préliminaires, le nombre d'appels augmente de manière exponentielle.
	
	En somme, la méthode des essais successifs est idéale pour se faire une idée du résultat avec des petites valeurs, puisqu'elle demande peu d'efforts de conception grâce à une trame fixe pour l'algorithme.
	En revanche, le nombre d'appels est exponentiel selon le nombre de tâches à ordonnancer.
	Il devient donc très long de faire tourner l'algorithme avec des grandes valeurs.

\section{Programmation dynamique}
	\paragraph{}
		On cherche le coût associé à l'affectation optimale de \emph{d} unités de temps aux tâches consécutives $T_{1} à T_{j}$. Soit \emph{coût-opt(j, d)} la fonction réalisant cette association.

	\subsection{Formule de récurrence}
		On s'appuie sur le \emph{principe d'optimalité de Bellman} : toute sous-solution d'une solution optimale est optimale. Ainsi, le coût de \emph{j} tâches auxquelles on accorde une durée totale \emph{d} de manière optimale est le minimum de la somme entre le coût de \emph{j-1} tâches auxquelles on accorde \emph{D-l} unités de temps, où \emph{l} parcourt l'ensemble des durées possibles pour la tâche \emph{j}, et le coût de la tâche \emph{j}.
	\subsection{Structure tabulaire}
	\subsection{Algorithme}
	\subsection{Complexité temporelle}
	\subsection{Complexité spatiale}

\section{Compléments}
\begin{enumerate}
\item
	Les algorithmes à essais successifs sont plus faciles à concevoir car ils suivent une trame fixe.
	Il suffit de déterminer les constituants de l'algorithme générique.
	En revanche, pour la méthode de programmation dynamique, il faut trouver une relation de récurrence qui n'est pas toujours évidente.
	Il est vrai qu'une fois qu'on l'a déterminé il devient aisé de produire l'algorithme mais cette étape est souvent dynamique.
	C'est pourquoi la méthode des essais successifs demande moins d'efforts de conception que la méthode de programmation dynamique qui est cependant plus efficace en terme de complexité temporelle.
\item
	D'après les résultats obtenus dans la phase d'expérimentation des algorithmes à essais successifs, on remarque qu'on peut aller en pratique jusqu'à une valeur $n = 50$ en fixant $k = 10$.
	On pouvait s'y attendre car la complexité au pire en termes de nombre de candidats à examiner est en $k^n$.
	Le nombre d'appels augmente donc de façon exponentielle, ce qui implique d'avoir une valeur de n faible pour pouvoir exécuter l'algorithme avec un temps raisonnable.
\item
	D'après les résultats obtenus dans la phase d'expérimentation des algorithmes à essais successifs, on remarque qu'on peut aller en pratique jusqu'à une valeur $k = 5000$ en fixant $n = 10$.
	Il est important de noter que la version sans élagage atteint une asymptote : cela est dû au fait que lorsque $k \ge D$, choisir une durée supérieure à D pour une tâche ne peut pas être satisfaisant.
	De plus, le nombre d'appels ne dépend pas des valeurs du tableau \emph{cd} pour la version sans élagage, ce qui n'est pas le cas de la version avec élagage.
	Cela peut s'expliquer par le rang des durées max possible qui varie selon les tirages aléatoires.
\item
\item
	Les algorithmes de type "diviser pour régner" sont basés sur une équation de récurrence de même type que ceux utilisant la programmation dynamique.
	Il est donc possible de résoudre le problème posé à l'aide d'une solution de type "diviser pour régner" en théorie.
	Cependant, la complexité spatiale serait très importante à cause du très grand nombre d'appels récursifs.
\end{enumerate}

\section{Conclusion}

En somme, nous avons pu expérimenter les différents aspects qui font d'un algorithme un bon algorithme.
En effet, nos solutions ont des complexités très différentes, mais le meilleur algorithme nécessite un effort de conception plus important.
Il est important de mettre en concurrence ces deux aspects car selon les données à traiter, il n'est pas forcément nécessaire d'avoir le meilleur algorithme qui soit.
Nous avons pu voir que sur de petites valeurs, les deux méthodes sont équivalentes en temps d'exécution.


\end{document}
